\chapter{Variables aleatorias continuas}\label{ch:g12:contdist}

Tiempos de espera, medidas físicas, proporciones: muchas cantidades
aleatorias toman un continuo de valores, y ningún valor aislado tiene
\omterm{def:g10:proba:distribution}{probabilidad} positiva. Las \omterm{def:g10:proba:distribution}{probabilidades} se calculan entonces integrando
una \emph{\omterm{def:g12:contdist:density}{densidad}}. Este capítulo estudia la \omterm{def:g12:contdist:uniform}{distribución uniforme}, la
exponencial y la \omterm{def:g12:contdist:normal}{distribución normal}, y usa esta última para cuantificar
las fluctuaciones de las encuestas y de las \omterm{def:g10:proba:sample}{muestras}.

\section{Densidades de probabilidad}

\begin{definition}[Densidad, variable aleatoria continua]\label{def:g12:contdist:density}
Una \emph{densidad de probabilidad}\index{densidad} sobre un \omterm{def:g10:numbers:interval}{intervalo} $I$
es una \omterm{def:g11:func:function}{función} $f$ \omterm{def:g12:limcont:continuity}{continua} y no negativa sobre $I$ con
$\int_I f(t)\,\dd t = 1$ (la \omterm{def:g12:integ:area}{integral} sobre un $I$ no acotado se entiende
como límite de \omterm{def:g12:integ:area}{integrales} sobre \omterm{def:g10:numbers:interval}{intervalos} acotados cada vez mayores). Una
\omterm{def:g12:randvar:rv}{variable aleatoria} $X$ tiene densidad $f$ si para todos $a \leq b$ de $I$:
\[
\P(a \leq X \leq b) = \int_a^b f(t)\,\dd t .
\]
\end{definition}

Las \omterm{def:g10:proba:distribution}{probabilidades} son \emph{áreas bajo la curva de \omterm{def:g12:contdist:density}{densidad}}. En
particular, $\P(X = a) = \int_a^a f = 0$ para cada valor aislado $a$: solo
los \omterm{def:g10:numbers:interval}{intervalos} llevan \omterm{def:g10:proba:distribution}{probabilidad}, y
$\P(a \leq X \leq b) = \P(a < X < b)$.

\begin{definition}[Esperanza y varianza]\label{def:g12:contdist:exp}
Para $X$ con \omterm{def:g12:contdist:density}{densidad} $f$ sobre $I$:
\[
\E(X) = \int_I t\,f(t)\,\dd t, \qquad
\V(X) = \int_I \bigl(t - \E(X)\bigr)^2 f(t)\,\dd t
= \E(X^2) - \E(X)^2 .
\]
\end{definition}

Las reglas de los capítulos 14 y 15 (linealidad de la \omterm{def:g12:randvar:exp}{esperanza},
aditividad de la \omterm{def:g12:randvar:exp}{varianza} para variables \omterm{def:g12:sums:indep}{independientes},
Bienaymé--Chebyshev) siguen siendo válidas; sus demostraciones, con
\omterm{def:g12:integ:area}{integrales} en lugar de sumas, se admiten a este nivel.

\section{Distribución uniforme}

\begin{definition}[Distribución uniforme]\label{def:g12:contdist:uniform}
$X$ sigue la \emph{distribución uniforme}\index{distribución uniforme}
$\mathcal U\intcc ab$ si su \omterm{def:g12:contdist:density}{densidad} es constante,
$f(t) = \frac{1}{b-a}$ sobre $\intcc{a}{b}$. Entonces, para
$\intcc{c}{d} \subseteq \intcc ab$,
$\P(c \leq X \leq d) = \frac{d - c}{b - a}$: \omterm{def:g10:proba:distribution}{probabilidad} proporcional a la
longitud.
\end{definition}

\begin{proposition}\label{prop:g12:contdist:uniformmoments}
Si $X \sim \mathcal U\intcc ab$, entonces
$\E(X) = \dfrac{a+b}{2}$ y $\V(X) = \dfrac{(b-a)^2}{12}$.
\end{proposition}

\begin{proof}
$\E(X) = \int_a^b \frac{t}{b-a}\dd t
= \frac{b^2 - a^2}{2(b-a)} = \frac{a+b}{2}$. Para la \omterm{def:g12:randvar:exp}{varianza},
$\E(X^2) = \frac{b^3 - a^3}{3(b-a)} = \frac{a^2 + ab + b^2}{3}$, y
\[
\V(X) = \frac{a^2 + ab + b^2}{3} - \frac{(a+b)^2}{4}
= \frac{4a^2 + 4ab + 4b^2 - 3a^2 - 6ab - 3b^2}{12}
= \frac{(b - a)^2}{12}. \qedhere
\]
\end{proof}

\section{Distribución exponencial}

\begin{definition}[Distribución exponencial]\label{def:g12:contdist:expo}
Para $\lambda > 0$, $X$ sigue la \emph{distribución
exponencial}\index{distribución exponencial} $\mathcal E(\lambda)$ si su
\omterm{def:g12:contdist:density}{densidad} sobre $\intco{0}{+\infty}$ es
\[
f(t) = \lambda\,\eu^{-\lambda t}.
\]
Entonces $\P(X \leq x) = 1 - \eu^{-\lambda x}$ y
$\P(X > x) = \eu^{-\lambda x}$ para $x \geq 0$.
\end{definition}

\begin{proposition}\label{prop:g12:contdist:expomoments}
Si $X \sim \mathcal E(\lambda)$:
$\E(X) = \dfrac{1}{\lambda}$ y $\V(X) = \dfrac{1}{\lambda^2}$.
\end{proposition}

\begin{proof}
Integrando por partes sobre $\intcc{0}{A}$:
\[
\int_0^A t\,\lambda\eu^{-\lambda t}\dd t
= \bigl[-t\,\eu^{-\lambda t}\bigr]_0^A + \int_0^A \eu^{-\lambda t}\dd t
= -A\eu^{-\lambda A} + \frac{1 - \eu^{-\lambda A}}{\lambda}
\xrightarrow[A\to+\infty]{} \frac1\lambda,
\]
usando $A\eu^{-\lambda A} \to 0$ (\cref{thm:g12:exp:growth}). Una segunda
\omterm{thm:g12:integ:ibp}{integración por partes} da $\E(X^2) = \frac{2}{\lambda^2}$, de donde
$\V(X) = \frac{2}{\lambda^2} - \frac{1}{\lambda^2}$.
\end{proof}

\begin{theorem}[Falta de memoria]\label{thm:g12:contdist:memoryless}
\index{falta de memoria}
Si $X \sim \mathcal E(\lambda)$, entonces, para todos $s, t \geq 0$:
\[
\pcond{X > s}{X > s + t} = \P(X > t) .
\]
La \omterm{def:g12:contdist:expo}{distribución exponencial} es la ley de las duraciones de vida \emph{sin
envejecimiento} (núcleos radiactivos, no bombillas).
\end{theorem}

\begin{proof}
\[
\pcond{X > s}{X > s+t}
= \frac{\P(X > s + t)}{\P(X > s)}
= \frac{\eu^{-\lambda(s+t)}}{\eu^{-\lambda s}} = \eu^{-\lambda t}
= \P(X > t). \qedhere
\]
\end{proof}

\begin{omfigure}
\begin{tikzpicture}
\begin{axis}[omaxis,
  width=8.8cm, height=5cm,
  xmin=-0.4, xmax=4.4, ymin=0, ymax=1.15,
  xtick={1.3}, xticklabels={$t$}, ytick={1}, yticklabels={$\lambda$}]
  \addplot[name path=dens, omDef, thick, domain=0:4.2, samples=80]
    {exp(-x)};
  \path[name path=xaxis] (0,0) -- (4.2,0);
  \addplot[omThm!20] fill between[of=dens and xaxis,
    soft clip={domain=1.3:4.2}];
  \draw[black, dashed, thin] (axis cs:1.3,0) -- (axis cs:1.3,0.2725);
  \node[omThm, font=\small, anchor=west] at (axis cs:2.1,0.35)
    {$\P(X > t) = \eu^{-\lambda t}$};
\end{axis}
\end{tikzpicture}

{\small La \omterm{def:g12:contdist:density}{densidad} exponencial $\lambda\eu^{-\lambda x}$ (aquí
$\lambda = 1$): el área de la cola más allá de $t$ (en rojo) vale
$\eu^{-\lambda t}$, y la \omterm{thm:g12:contdist:memoryless}{falta de memoria} dice que toda cola se parece a la
\omterm{def:g12:randvar:rv}{distribución} entera reescalada.}
\end{omfigure}

\section{Distribución normal}

\begin{definition}[Distribución normal]\label{def:g12:contdist:normal}
$X$ sigue la \emph{distribución normal}\index{distribución normal}
\emph{estándar} $\mathcal N(0, 1)$ si su \omterm{def:g12:contdist:density}{densidad} sobre $\R$ es
\[
\varphi(t) = \frac{1}{\sqrt{2\pi}}\,\eu^{-t^2/2}
\]
(la campana). Más en general, $X \sim \mathcal N(\mu, \sigma^2)$ si
$\dfrac{X - \mu}{\sigma} \sim \mathcal N(0,1)$; entonces $\E(X) = \mu$ y
$\V(X) = \sigma^2$.
\end{definition}

\begin{omfigure}
\begin{tikzpicture}
\begin{axis}[
  width=11cm, height=4.6cm,
  axis lines=middle, xmin=-4, xmax=4, ymin=0, ymax=0.45,
  xtick={-3,-2,-1,0,1,2,3},
  xticklabels={$\mu{-}3\sigma$,$\mu{-}2\sigma$,$\mu{-}\sigma$,$\mu$,
               $\mu{+}\sigma$,$\mu{+}2\sigma$,$\mu{+}3\sigma$},
  ytick=\empty, samples=120, domain=-4:4]
  \addplot[name path=gauss, omDef, thick] {exp(-x^2/2)/sqrt(2*pi)};
  \path[name path=xaxis] (-4,0) -- (4,0);
  \addplot[omDef!20] fill between[of=gauss and xaxis,
           soft clip={domain=-2:2}];
  \node[omDef] at (axis cs:0,0.17) {$\approx 95.4\,\%$};
\end{axis}
\end{tikzpicture}

{\small La \omterm{def:g12:contdist:density}{densidad} normal: en torno al $68.3\,\%$ de la masa está a menos
de $\sigma$ de la \omterm{def:g11:stat:mean}{media}, el $95.4\,\%$ a menos de $2\sigma$ y el
$99.7\,\%$ a menos de $3\sigma$.}
\end{omfigure}

\begin{remark}
El factor $\frac{1}{\sqrt{2\pi}}$ hace que el área total valga $1$: un
cálculo famoso (la \emph{\omterm{def:g12:integ:area}{integral} de Gauss}) que se hace en la universidad.
No hay fórmula elemental para $\int \varphi$: las \omterm{def:g10:proba:distribution}{probabilidades} normales
se leen en tablas o en la calculadora.
\end{remark}

\begin{theorem}[De Moivre--Laplace, admitido]\label{thm:g12:contdist:tcl}
\index{teorema de De Moivre--Laplace}
Sean $X_n \sim \mathcal B(n, p)$ y
$Z_n = \dfrac{X_n - np}{\sqrt{np(1-p)}}$ (la binomial tipificada).
Entonces, para todos $a \leq b$,
\[
\P(a \leq Z_n \leq b) \xrightarrow[n\to+\infty]{}
\int_a^b \varphi(t)\,\dd t .
\]
\emph{Este resultado se admite a este nivel.} Explica la aparición
universal de la campana: una binomial con $n$ grande es aproximadamente
normal, y (teorema central del límite, en la universidad) también lo es
cualquier suma de muchos efectos pequeños e \omterm{def:g12:sums:indep}{independientes}.
\end{theorem}

\begin{method}[Intervalo de fluctuación al $95\,\%$]\label{met:g12:contdist:fluctuation}
\index{intervalo de confianza}
Como una variable normal cae a menos de $1.96$ desviaciones típicas de su
\omterm{def:g11:stat:mean}{media} con \omterm{def:g10:proba:distribution}{probabilidad} $0.95$, para $n$ grande una frecuencia
$F_n = \frac{X_n}{n}$ de una \omterm{def:g10:proba:sample}{muestra} $\mathcal B(n,p)$ cumple
\[
\P\left(p - 1.96\sqrt{\tfrac{p(1-p)}{n}} \leq F_n \leq
p + 1.96\sqrt{\tfrac{p(1-p)}{n}}\right) \approx 0.95 .
\]
\begin{itemize}
  \item \emph{\omterm{def:g10:numbers:interval}{Intervalo} de fluctuación} (para contrastar): si un $p$
        afirmado deja la frecuencia observada fuera de este \omterm{def:g10:numbers:interval}{intervalo}, se
        rechaza la afirmación al nivel del $5\,\%$.
  \item \emph{\omterm{met:g12:contdist:fluctuation}{Intervalo de confianza}} (para estimar): el \omterm{def:g10:numbers:interval}{intervalo}
        simplificado
        $\left[F_n - \frac{1}{\sqrt n},\ F_n + \frac{1}{\sqrt n}\right]$
        contiene $p$ con \omterm{def:g10:proba:distribution}{probabilidad} al menos $0.95$ (usando
        $1.96\sqrt{p(1-p)} \leq 1$; véase el \cref{exo:g12:sums:8}).
\end{itemize}
\end{method}

\begin{example}\label{ex:g12:contdist:poll}
Una encuesta a $n = 1000$ votantes da a un candidato un $52\,\%$. El
\omterm{met:g12:contdist:fluctuation}{intervalo de confianza}
$52\,\% \pm \frac{1}{\sqrt{1000}} \approx 52\,\% \pm 3.2\,\%$ contiene
todavía el $50\,\%$: la encuesta, por sí sola, no demuestra que el
candidato vaya por delante.
\end{example}

\section{Ejercicios}

\begin{exercise}[$\star$]\label{exo:g12:contdist:1}
Un autobús pasa cada $15$ minutos; una viajera llega en un instante
elegido al azar de manera uniforme. Sea $X \sim \mathcal U\intcc{0}{15}$ el
tiempo de espera. Calcula $\P(X \leq 5)$, $\P(4 \leq X \leq 10)$, $\E(X)$ y
$\sigma(X)$.
\end{exercise}

\begin{exercise}[$\star$]\label{exo:g12:contdist:2}
Comprueba que $f(t) = \frac{3}{8}t^2$ es una \omterm{def:g12:contdist:density}{densidad} sobre $\intcc{0}{2}$
y calcula $\E(X)$ y $\P(X \geq 1)$ para una variable $X$ con esa \omterm{def:g12:contdist:density}{densidad}.
\end{exercise}

\begin{exercise}[$\star$]\label{exo:g12:contdist:3}
La duración de vida (en años) de un componente electrónico sigue
$\mathcal E(0.2)$.
\begin{enumerate}
  \item Calcula la duración de vida esperada y $\P(X > 5)$.
  \item El componente lleva ya $3$ años funcionando. ¿Cuál es la
        \omterm{def:g10:proba:distribution}{probabilidad} de que funcione al menos $5$ años más?
\end{enumerate}
\end{exercise}

\begin{exercise}[$\star\star$]\label{exo:g12:contdist:4}
El periodo de semidesintegración de un átomo radiactivo cuya duración de
vida sigue $\mathcal E(\lambda)$ es la \omterm{def:g11:stat:median}{mediana} $m$: $\P(X > m) = \frac12$.
Expresa $m$ en \omterm{def:g11:func:function}{función} de $\lambda$ y compárala con la \omterm{def:g12:randvar:exp}{esperanza}. ¿Cuál es
mayor y por qué tiene sentido en una \omterm{def:g12:randvar:rv}{distribución} asimétrica?
\end{exercise}

\begin{exercise}[$\star\star$]\label{exo:g12:contdist:5}
Las estaturas de una población siguen $\mathcal N(175,\ 7^2)$ (en cm).
Usando la regla $68$--$95$--$99.7$, estima la proporción de la población
con estatura entre $168$ y $182$ cm, por encima de $189$ cm y por debajo de
$154$ cm.
\end{exercise}

\begin{exercise}[$\star\star$]\label{exo:g12:contdist:6}
Una máquina llena bolsas etiquetadas de $500$ g; la masa envasada sigue
$\mathcal N(\mu,\ 4^2)$. La normativa exige que como mucho el $2.5\,\%$ de
las bolsas pese menos de $500$ g. Usando la regla del $95\,\%$, ¿cuál es el
ajuste \omterm{def:g10:functions:extrema}{mínimo} de $\mu$ que cumple la \omterm{def:g11:scal:dot}{norma}?
\end{exercise}

\begin{exercise}[$\star\star$]\label{exo:g12:contdist:7}
Un dado del que se afirma que está equilibrado se lanza $1200$ veces y saca
seis $260$ veces (como en el \cref{exo:g12:sums:7}).
\begin{enumerate}
  \item Calcula el \omterm{def:g10:numbers:interval}{intervalo} de fluctuación al $95\,\%$ de la frecuencia de
        seises de un dado equilibrado en $1200$ lanzamientos.
  \item ¿Está dentro la frecuencia observada? Compara la fuerza de esta
        conclusión con la del análisis mediante Bienaymé--Chebyshev.
\end{enumerate}
\end{exercise}

\begin{exercise}[$\star\star\star$]\label{exo:g12:contdist:8}
Antes de unas elecciones, una encuesta a $n$ personas va a estimar el
resultado $p$ de un candidato mediante la frecuencia observada $F_n$.
\begin{enumerate}
  \item Con el \omterm{met:g12:contdist:fluctuation}{intervalo de confianza} del
        \cref{met:g12:contdist:fluctuation}, ¿qué tamaño de \omterm{def:g10:proba:sample}{muestra}
        garantiza un margen de $\pm 2$ puntos?
  \item En realidad, a los candidatos los separa $1$ punto. Explica por qué
        ninguna encuesta realista puede declarar con fiabilidad al ganador,
        por bien hecha que esté.
\end{enumerate}
\end{exercise}

\begin{exercise}[$\star\star\star$]\label{exo:g12:contdist:9}
Sea $X$ de \omterm{def:g12:contdist:density}{densidad} $f$ sobre $\intcc ab$ y sea $Y = \alpha X + \beta$ con
$\alpha > 0$.
\begin{enumerate}
  \item Demuestra que $\P(c \leq Y \leq d)
        = \P\left(\frac{c - \beta}{\alpha} \leq X \leq
        \frac{d-\beta}{\alpha}\right)$, y deduce que $Y$ tiene \omterm{def:g12:contdist:density}{densidad}
        $g(y) = \frac{1}{\alpha} f\left(\frac{y - \beta}{\alpha}\right)$.
  \item Deduce que si $Z \sim \mathcal N(0,1)$, entonces
        $\mu + \sigma Z$ tiene \omterm{def:g12:contdist:density}{densidad}
        $\frac{1}{\sigma\sqrt{2\pi}}\,
        \eu^{-(y - \mu)^2/(2\sigma^2)}$: la \omterm{def:g12:contdist:density}{densidad} normal general.
\end{enumerate}
\end{exercise}

\section{Problema: Esperar el autobús, medir a la multitud}

\begin{problem}[{Problema de fin de semana --- tres densidades gobiernan
el mundo: la ignorancia total, la espera sin memoria y la campana de las
muchas causas pequeñas; más la paradoja por la que tu autobús siempre llega
tarde}]
\label{pb:g12:contdist:1}
¿Por qué parece que tu autobús tarda siempre más de lo que promete el
horario? ¿Por qué la clase en la que se sienta un estudiante medio es más
grande que la clase \omterm{def:g11:stat:mean}{media}? ¿Por qué tus amigos tienen más amigos que tú?
Una pieza astuta de matemáticas, el muestreo sesgado por la longitud,
responde a las tres, y vive en este capítulo, entre la \omterm{def:g12:contdist:uniform}{distribución
uniforme}, la exponencial y la \omterm{def:g12:contdist:normal}{distribución normal} (\cref{def:g12:contdist:uniform},
\cref{def:g12:contdist:expo}, \cref{def:g12:contdist:normal}). Este último
problema del volumen pone a trabajar a las tres \omterm{def:g12:contdist:density}{densidades} y termina donde
apunta toda la serie: la campana.

\medskip
\noindent\textbf{Parte I --- Tres personajes.}
\begin{enumerate}
  \item Un autobús llega en un instante elegido al azar de manera uniforme
        en los próximos $30$ minutos. Calcula
        $\P(\text{espera} \leq 10)$, la espera esperada y $\sigma$
        (\cref{prop:g12:contdist:uniformmoments}).
  \item Comprueba que $f(x) = 2x$ sobre $\intcc{0}{1}$ es una \omterm{def:g12:contdist:density}{densidad} y
        calcula $\E(X)$ y $\P\left(X > \frac12\right)$.
  \item Las llamadas telefónicas llegan a un teléfono de ayuda como un
        flujo sin memoria: la espera hasta la siguiente llamada es
        exponencial de \omterm{def:g11:stat:mean}{media} $10$ minutos. Calcula $\P(T > 15)$ y la espera
        \omterm{def:g11:stat:median}{mediana}; ¿por qué es la \omterm{def:g11:stat:median}{mediana} \emph{menor} que la \omterm{def:g11:stat:mean}{media}?
  \item Ya has esperado $20$ minutos. ¿Cuánto vale
        $\P(T > 35 \mid T > 20)$ (\cref{thm:g12:contdist:memoryless})?
        Interprétalo en una frase.
  \item Empareja cada espera con su modelo correcto (uniforme, exponencial
        o ninguno de los dos): (a) el chasquido de un contador Geiger; (b)
        un metro que pasa cada $8$ minutos, con tu llegada sin
        sincronizar; (c) la siguiente llamada al teléfono de ayuda; (d) el
        fallo de una bombilla que \emph{se desgasta}. Justifica el caso (d)
        con el teorema de la \omterm{thm:g12:contdist:memoryless}{falta de memoria}.
\end{enumerate}

\medskip
\noindent\textbf{Parte II --- La campana en acción.}
(Usa las referencias de la normal: $68\,\%$ a menos de $\sigma$,
$95\,\%$ a menos de $2\sigma$ y $99.7\,\%$ a menos de $3\sigma$.)
\begin{enumerate}[resume]
  \item Las estaturas adultas siguen aproximadamente
        $\mathcal N(172,\ 8)$ (cm). ¿Qué fracciones caen en
        $\intcc{164}{180}$, $\intcc{156}{188}$ y $\intcc{148}{196}$?
  \item Estima la proporción de personas más altas de $190$ cm (primero la
        \omterm{pb:g11:stat:1}{puntuación z}; después, tabla o calculadora).
  \item Las puntuaciones de cociente intelectual se calibran a
        $\mathcal N(100, 15)$. ¿Qué fracción supera $130$? ¿Y
        aproximadamente una persona de cada cuántas?
  \item Una máquina corta tornillos con longitudes
        $\mathcal N(50,\ 0.1)$ mm y la especificación es
        $\intcc{49.8}{50.2}$. ¿Qué proporción se rechaza? La industria
        celebra los procesos ``seis sigma'', cuyas especificaciones están a
        $\pm 6\sigma$: ¿qué se compra con eso y por qué lo persiguen las
        fábricas?
  \item De Moivre--Laplace (\cref{thm:g12:contdist:tcl}): para $100$
        lanzamientos de una moneda equilibrada, aproxima
        $\P(45 \leq \text{caras} \leq 55)$ (corrección por \omterm{def:g12:limcont:continuity}{continuidad}:
        $\pm 0.5$; $\sigma = 5$). ¿Qué máquina de madera del
        \cref{pb:g11:binom:1} convierte este teorema en una curva suave?
  \item Del \cref{exo:g12:contdist:8}: un margen de $\pm 2$ puntos
        necesita unas $n = 2\,400$ personas. Calcula la \omterm{def:g10:proba:sample}{muestra} necesaria
        para $\pm 0.5$ puntos y concluye, en una frase, por qué ninguna
        encuesta puede declarar honradamente el ganador de una carrera
        separada por un punto.
\end{enumerate}

\medskip
\noindent\textbf{Parte III --- La paradoja del autobús.}
\begin{enumerate}[resume]
  \item Los autobuses llegan como un flujo sin memoria con un \omterm{def:g10:numbers:interval}{intervalo}
        medio de $10$ minutos. Tú llegas a la parada en un momento
        cualquiera. La intuición dice que la espera esperada es de $5$
        minutos (medio \omterm{def:g10:numbers:interval}{intervalo}). ¿Qué dice en cambio el teorema de la
        \omterm{thm:g12:contdist:memoryless}{falta de memoria}?
  \item Diagnostica la intuición con un horario de juguete: los \omterm{def:g10:numbers:interval}{intervalos}
        alternan $5$ y $15$ minutos (\omterm{def:g10:numbers:interval}{intervalo} medio, $10$). Calcula la
        \omterm{def:g10:proba:distribution}{probabilidad} de que una llegada uniforme al azar caiga en un
        \omterm{def:g10:numbers:interval}{intervalo} largo y después la espera esperada de verdad; y nombra al
        culpable: las llegadas muestrean los \omterm{def:g10:numbers:interval}{intervalos} con \omterm{def:g10:proba:distribution}{probabilidad}
        proporcional a su \emph{longitud}.
  \item En el flujo sin memoria, el \omterm{def:g10:numbers:interval}{intervalo} en el que caes tiene una
        longitud esperada de $20$ minutos, el doble del \omterm{def:g10:numbers:interval}{intervalo} típico
        (el tiempo hacia atrás hasta el autobús anterior y el tiempo hacia
        delante hasta el siguiente son \emph{los dos} exponenciales de
        \omterm{def:g11:stat:mean}{media} $10$). Concilia esto con la pregunta 12 y enuncia en una
        frase la \emph{paradoja de la inspección}.
  \item El mismo aguijón, en tierra firme: una universidad imparte nueve
        clases de $10$ estudiantes y una de $110$. Calcula el tamaño medio
        de clase y después el tamaño de clase que vive el estudiante
        \emph{medio}. ¿Qué número imprimirá el folleto y cuál es la verdad
        vivida por los estudiantes?
  \item La paradoja de la amistad (``tus amigos tienen, en promedio, más
        amigos que tú'') son las mismas matemáticas. Di en una frase qué
        hace aquí el papel del \omterm{def:g10:numbers:interval}{intervalo} largo.
\end{enumerate}

\medskip
\noindent\textbf{Parte IV --- Cerrando el volumen.}
\begin{enumerate}[resume]
  \item La simulación, el puente de quienes aplican esto: si $U$ es
        uniforme sobre $\intoo{0}{1}$, demuestra que $T = -10\ln U$ es
        exponencial de \omterm{def:g11:stat:mean}{media} $10$ (calcula $\P(T > t)$). Todas las
        simulaciones industriales de esperas aleatorias funcionan con este
        truco de una línea.
  \item Una vieja receta de programadores construye una campana: se suman
        doce uniformes \omterm{def:g12:sums:indep}{independientes} sobre $\intcc{0}{1}$ y se resta $6$.
        Da la \omterm{def:g11:stat:mean}{media} y la \omterm{def:g12:randvar:exp}{varianza} de la suma y explica, con la tabla de
        Galton en la mano, por qué el resultado es casi normal.
  \item Los mercados se desploman con más violencia de la que permite la
        campana: el desplome de 1987 se describió como un ``\omterm{def:g11:prob:model}{suceso} de
        $25\sigma$'', que bajo normalidad tiene una \omterm{def:g10:proba:distribution}{probabilidad} del orden
        de $10^{-137}$. ¿Cuál es la conclusión correcta: sobre el mundo o
        sobre el modelo? (Una frase: el catecismo del estadístico.)
  \item Final y despedida del volumen: los tres personajes, a una línea
        cada uno: la uniforme (ignorancia dentro de unos límites), la
        exponencial (un riesgo sin memoria) y la normal (la suma de muchas
        causas pequeñas); la paradoja de la inspección como aguijón del
        capítulo; y el arco de toda la serie, desde contar cartones de
        zumo en el año 6 hasta la campana que mide a cualquier multitud,
        con los volúmenes universitarios esperando allí donde toman el
        relevo los límites, las \omterm{def:g12:integ:area}{integrales} y el teorema central del
        límite.
\end{enumerate}
\end{problem}
